\section{Способы определения коэффициента оптического поглощения нелинейно-оптических кристаллов}\label{sec:review/abs}

При разработке и исследованиях мощных лазерных источников, работа которых основана на нелинейно-оптическом преобразовании света, в целом и источников среднего ИК-диапазона в частности, крайне важную роль играют эффекты, связанные с тепловым воздействием и самовоздействием на нелинейно-оптический кристалл. В условиях высокой мощности лазерного излучения даже невысокий коэффициент оптического поглощения (а для кристалла PPLN на длине волны 3~мкм он составляет $10^{-2}~\text{см}^{-1}$ \cite{ostapiv2021determination}) приводит к проявлению термических эффектов. К ним следует отнести тепловую линзу \cite{peng2015high}, и, что более критично на невысоких уровнях оптической мощности,~--- изменение показателя преломления и нарушение условий фазового синхронизма, влекущее за собой снижение эффективности нелинейно-оптического преобразования и ухудшение качества выходного пучка. Более того, указанный эффект обладает обратной связью по мощности преобразованного излучения, что при некоторых условиях может приводить к бистабильности системы и хаотическому поведению эффективности преобразования \cite{dmitriev2004prikladnaya, zotov2024equivalent}. Таким образом, точное измерение коэффициента поглощения позволяет оптимизировать режимы работы лазерной системы, выбрать оптимальные кристаллы и предотвратить их термическое повреждение, обеспечивая стабильность и надежность мощных лазерных источников среднего ИК-диапазона.




\subsection{Лазерная калориметрия}

\subsection{Пьезорезонансная лазерная калориметрия и её модификации}
	\nomenclature{ПРЛК}{Пьезорезонансная лазерная калориметрия}
	

